\chapter{Cronograma de ejecucion}
\thispagestyle{memoria}\pagestyle{memoria}

\section{Plazo referencial}
Para el alcance electrico del anteproyecto se adopta un plazo referencial de ocho
semanas. Es una secuencia academica y debe ajustarse al cronograma civil, a la
disponibilidad de equipos certificados y a la autorizacion para intervenir cerca
de combustibles.

\begin{table}[H]
\centering
\small
\begin{tabularx}{\textwidth}{@{}L*{8}{c}@{}}
\toprule
Actividad & 1 & 2 & 3 & 4 & 5 & 6 & 7 & 8 \\
\midrule
Compatibilizacion, permisos y replanteo & X & X & & & & & & \\
Canalizaciones enterradas y registros & & X & X & X & & & & \\
Canalizaciones del edificio administrativo & & X & X & X & & & & \\
Anillo PAT, electrodos y equipotencialidad & & & X & X & X & & & \\
Tendido de alimentadores y circuitos & & & & X & X & X & & \\
Montaje de tableros, luminarias y dispositivos & & & & & X & X & X & \\
Grupo, ATS, UPS y paro de emergencia & & & & & & X & X & \\
Pruebas, rotulado, correcciones y entrega & & & & & & & X & X \\
\bottomrule
\end{tabularx}
\caption{Cronograma referencial de ocho semanas.}
\end{table}

\section{Hitos de control}
Antes de tapar zanjas se verifican profundidad, separaciones, continuidad de
ductos y registros. Antes del tendido se prueba cada canalizacion con guia. Antes
de energizar se aprueban aislamiento, continuidad del PE, resistencia de PAT,
secuencia de fases, disparo diferencial, paro de emergencia y transferencia. La
recepcion final incluye planos conforme a obra y protocolos.
